% Part: incompleteness
% Chapter: representability-in-q
% Section: beta-function

\documentclass[../../../include/open-logic-section]{subfiles}

\begin{document}

\olfileid{inc}{req}{bet}
\olsection{The Beta Function Lemma}

为了证明加法、乘法和 $\Char{=}$ 足以实现原始递归，我们需要构造能够处理序列的函数。（如果我们还有指数运算，任务会简单些。）当我们有原始递归时，可以定义诸如“第 $n$ 个质数”之类的概念，并采用相当直接的编码方式。但在这里我们没有原始递归——事实上我们正是要证明利用极小化即可实现原始递归——因此需要更巧妙的方法。

\begin{lem}
\ollabel{lem:beta}
存在一个函数 $\beta(d,i)$，使得对每个序列 $a_0$,
\dots,~$a_n$，都存在一个数~$d$，满足对每个 $i \le n$，
$\beta(d,i) = a_i$。并且，$\beta$ 可以仅由基本函数通过复合与正则极小化来定义。
\end{lem}

不妨将 $d$ 看作序列 $\tuple{a_0, \dots, a_n}$ 的编码，
而 $\beta(d,i)$ 返回第 $i$ 个元素。（注意，此处的“编码”\emph{并非}使用我们已经熟悉的素数幂编码！）
这个引理的要求是极弱的；它并没有要求我们可以对序列进行拼接或追加，甚至没有要求我们能利用复合与正则极小化可定义的函数从 $a_0$, \dots,~$a_n$ \emph{计算}出~$d$。
它只断言存在一个“解码”函数，使得每个序列都有其“编码”。

记号 $\beta$ 的使用源于 G\"odel。再次强调，证明此引理的困难之处在于，要用看似受限的资源（即仅用复合和极小化——尽管我们可以使用加法、乘法和~$\Char{=}$）定义出合适的~$\beta$。
有多种方法可以证明这个引理，但最简洁的方法之一仍然是 G\"odel 的原始方法，它用到了一个称为孙子定理（传统上称为“中国剩余定理”）的数论事实。

\begin{defn}
两个自然数 $a$ 和 $b$ 是\emph{互素}的，当且仅当它们的最大公约数是~$1$；换言之，它们没有其他公共的除数。
\end{defn}

\begin{defn}
自然数 $a$ 与 $b$ \emph{模 $c$ 同余}，记作 $a \equiv b \mod c$，当且仅当 $c \mid (a-b)$，即 $a$ 和 $b$ 除以 $c$ 的余数相同。
\end{defn}

以下是孙子定理：

\begin{thm}
设 $x_0$，\dots，$x_n$（两两）互素。令 $y_0$，\dots，$y_n$ 为任意数。则存在一个数 $z$，使得
\begin{align*}
z & \equiv y_0 \mod x_0 \\
z & \equiv y_1 \mod x_1 \\
& \vdots  \\
z & \equiv y_n \mod x_n.
\end{align*}
\end{thm}

下面说明我们如何使用孙子定理：如果 $x_0$，\dots，$x_n$ 分别大于 $y_0$，\dots，$y_n$，那么我们可以取 $z$ 来编码序列 $\tuple{y_0, \dots,y_n}$。要恢复 $y_i$，只需将 $z$ 除以 $x_i$ 并取余数。为了使用这种编码，我们需要为 $x_0$，\dots，$x_n$ 找到合适的取值。

以下几个观察将对我们有所帮助。给定 $y_0$，\dots，$y_n$，令
\begin{align*}
j &= \max(n, y_0 + 1, \dots, y_n + 1), \\
m &= \lcm(1,\dots,j),
\end{align*}
再令
\begin{align*}
x_0 & = 1 + m \\
x_1 & = 1 + 2 \cdot m \\
x_2 & = 1 + 3 \cdot m \\
& \vdots  \\
x_n & = 1 + (n+1) \cdot m
\end{align*}
则以下两点成立：

\begin{enumerate}
\item\ollabel{rel-prime} $x_0,\dots,x_n$ 是互素的。
\item\ollabel{less} 对每个 $i$，$y_i < x_i$。
\end{enumerate}

要说明 \olref{rel-prime} 成立，注意到：若素数 $p$ 满足 $p \mid x_i$ 与 $p \mid x_k$，则 $p \mid 1 + (i+1) m$ 且 $p \mid 1 + (k+1) m$。于是 $p$ 整除两者的差，\[
(1 + (i+1)m) - (1+ (k+1)m) = (i-k) m.
\]。由于 $p$ 整除 $1 + (i+1)m$，它必不能同时整除 $m$（否则前一次除法会留下余数 $1$）。既然 $p$ 整除 $(i-k)m$，所以 $p$ 整除 $i-k$。但是 $\left|i-k\right|$ 至多是 $n$，而我们已选取 $j \geq n$，这意味着 $p \mid m$，再次导致矛盾。因此没有素数能同时整除 $x_i$ 和 $x_k$。\olref{less} 是显然的：我们有 $y_i < j \leq m < x_i$。

现在我们来证明 $\beta$ 函数引理。注意我们可以使用 $0$、后继、加法、乘法、$\Char{=}$、投影函数，以及任何通过复合和对正则函数应用极小化而定义的函数。如果一个关系的特征函数可以这样定义，我们也可以使用该关系。如前所述，我们可以证明这些关系在布尔组合和有界量词下是封闭的；例如：
\begin{align*}
\fn{not}(x) & \defis \Char{=}(x,0)\\
\bmin{x \leq z}{R(x,y)} & \defis \umin{x}{(R(x,y) \lor x = z)}\\
\bexists{x \leq z}{R(x,y)} & \defiff R(\bmin{x \leq z}{R(x,y)}, y)
\end{align*}
进而我们可以证明，以下各项也都可以无需原始递归即可定义：

\begin{enumerate}
\item 配对函数，$J(x,y) = \frac{1}{2}[(x+y)(x+y+1)] + x$；
% maybe explain more what is going on here, a bit confusing.
\item 投影函数
\begin{align*}
K(z) & = \bmin{x \leq z}{\bexists{y \leq z}{z = J(x,y)}},\\
L(z) & = \bmin{y \leq z}{\bexists{x \leq z}{z = J(x,y)}};
\end{align*}
\item 小于关系 $x < y$；
\item 整除关系 $x \mid y$；
% \item $x \tsub y$
% \item $\fn{Prime}(x)$
% \item Assuming $p$ is prime, the relation ``$x$ is a power of $p$'':
% \[
% \bforall{y \leq x}{(y \mid x \lif y = 1 \lor y = x)}.
% \]
\item 函数 $\fn{rem}(x,y)$，它返回 $y$ 除以 $x$ 的余数。
\end{enumerate}

现在我们定义
\begin{align*}
\beta^*(d_0,d_1,i) & = \fn{rem}(1+(i+1) d_1,d_0) \text{ 且 }\\
\beta(d,i) & = \beta^*(K(d),L(d),i).
\end{align*}
这就是我们想要的函数。给定上述 $a_0,\dots,a_n$，令
\[
j = \max(n,a_0+1,\dots,a_n+1),
\]
并令 $d_1 = \lcm(1,\dots,j)$。由上文 \olref{rel-prime} 可知，
$1+d_1$， $1+2 d_1$，\dots， $1+(n+1) d_1$ 是互素的，
且由~\olref{less} 可知它们都大于 $a_0,\dots,a_n$。
根据孙子定理，存在一个值~$d_0$，使得对每个~$i$，
\[
d_0 \equiv a_i \mod (1+(i+1)d_1)
\]
于是（因为 $d_1$ 大于~$a_i$）有
\[
a_i = \fn{rem}(1+(i+1)d_1,d_0).
\]
令 $d = J(d_0,d_1)$。则对每个 $i \le n$，我们有
\begin{align*}
\beta(d,i) & =  \beta^*(d_0,d_1,i) \\
& =  \fn{rem}(1+(i+1) d_1,d_0) \\
& =  a_i
\end{align*}
这正是我们所需要的。至此，$\beta$ 函数引理的证明完成。

\begin{prob}
  证明关系 $x < y$、$x \mid y$ 以及函数~$\fn{rem}(x,y)$ 可以在不使用原始递归的情况下定义。
  你可以使用 $0$、后继、加法、乘法、$\Char{=}$、投影函数以及有界极小化和有界量词。
\end{prob}

\end{document}